\section{发布网络的候选局部化定理}\label{sec:theorem}
\subsection{可检查的结构条件}
线程构成无自环强连通有向网络 $\mathcal D$，线程数 $p\ge2$，有向直径为 $d$。边 $u\to t$ 表示线程 $t$ 读取 $u$ 的标志。每个非初始化事件有阶段号 $\sigma(e)\in\{1,\ldots,n\}$，线程内阶段依程序序递增。对所有线程给出完整的 $n$ 阶段输入。以下条件定义本文研究的发布网络。

\paragraph{P1：独占标志的重复发布。}
线程 $t$ 独占标志地址 $f_t$，初值 0；数据地址与所有标志地址分离。在阶段 $i$，$t$ 对 $f_t$ 恰有 $m(t,i)\ge2$ 次 release 写，值均为 $i\bmod2$。这些写按 $\SB$ 排列，覆盖该标志的全部非初始化写，没有阶段外隐藏写。发布之间允许固定的数据访问，不要求相邻，也不要求语义上可合并。

\paragraph{P2：完整交替观测。}
线程 $t$ 的阶段 $i$ 在最后一次发布之后，对每条入边 $u\to t$ 各有一次 acquire 读 $r(u,t,i)$，观测 $(i-1)\bmod2$。这些是标志地址的全部读；其来源未预先指定。等待之后仍可有数据访问。这里给出的是完整有限片段，而非缺少未来观测的在线前缀。

\paragraph{P3：阶段密度与写覆盖。}
每阶段所有线程合计至多 $b$ 个非初始化事件。每个数据地址的活动区间是其最早到最晚访问阶段；该区间中的每个完整 $G$ 阶段子区间至少含一次该地址的写，$G\ge1$。标志每阶段有写，自动满足此条件。初始化写数为 $a$。参数 $b,d,G,a$ 可以固定，$n$ 可增长。

条件不预设候选图小宽度，也不要求数据地址单写者。尤其，P3 限制的是写的时间覆盖密度，并不指定不同线程写之间的 CO 顺序。图~\ref{fig:proofchain} 概括各条件在证明中的作用。

\begin{figure}[tb]
\centering
\begin{tikzpicture}[node distance=4mm,box/.style={draw=black!50,rounded corners=2pt,align=center,text width=9.6cm,inner sep=5pt,font=\small},>=Latex]
\node[box,fill=blue!4] (a) {P1、P2：完整交替读链 $\Longrightarrow$ 来源块号锁定，块内身份保留};
\node[box,below=of a] (b) {$R_\forall$ 与网络直径 $d$ $\Longrightarrow$ 共同因果分隔 $h=2d+1$};
\node[box,below=of b] (c) {分隔与写覆盖 $G$ $\Longrightarrow$ RF、未决 CO、$\Cov(L)$ 跨度 $\le\delta$};
\node[box,below=of c,fill=blue!4] (d) {阶段密度 $b$ 与初始化数 $a$ $\Longrightarrow$ 全候选接口 $k=(\delta+1)b+a$};
\draw[->] (a)--(b);\draw[->] (b)--(c);\draw[->] (c)--(d);
\end{tikzpicture}
\caption{核心证明链。$\delta=2h+G$；最后再接精确编译，局部化本身不等于一致性判定。}\label{fig:proofchain}
\end{figure}

\subsection{定理陈述}
记 $N=|E|$，$I$ 为显式输入规模，包含事件、候选和链描述。候选表示仅物化 $D$ 中的 RF、$L$ 未决定的 CO 两向选择、$\Cov(L)$、相邻 $\SB$ 和稀疏初始化边；不把稠密传递闭包 $H,L$ 的每个关系对都物化为边。事件按阶段布局，选择组按最大端点的布局位置排序。组处理时仍可能参与后续组的事件，与本组端点一起组成事件袋。

\begin{cthm}[来源未定发布网络的局部化]\label{thm:main}
对满足 P1--P3 的输入，可靠核 $\ka$ 或证明 UNSAT，或认证下列性质：
\begin{enumerate}[label=(\Alph*)]
\item 每个 $r(u,t,i)$ 的来源域包含于 $u$ 的第 $i-1$ 阶段发布块；$i=1$ 时只含初始化写。若 $i\ge2$，该块的首次发布 $H$-先于 $r(u,t,i)$，即使域仍有多个来源。
\item 令 $h=2d+1$。若非初始化事件 $e,f$ 满足 $\sigma(f)-\sigma(e)\ge h$，则 $H(e,f)$。
\item 令 $\delta=2h+G=4d+2+G$。上述候选表示中所有非初始化边的阶段跨度至多 $\delta$。二元组布局的事件袋大小至多
\begin{equation}
 k=(\delta+1)b+a=(4d+G+3)b+a.\label{eq:k}
\end{equation}
\item 存在保留全部原始 RF／CO 支持的编译图与支持 CNF。置 $S=2^{33k^2+k}$，图节点、边、标签文字及 CNF 文字出现次数均可界为
\begin{equation}
 O(|X|+kNS),\label{eq:size}
\end{equation}
保守构造时间为 $O(I^5+|X|+kNS\operatorname{poly}(k))$。对任意 $\alpha\subseteq\Lit(X)$，支持扫描或支持 CNF 的纯 UP 检出无解，或给出全部必然原始文字。
\end{enumerate}
\end{cthm}

结论 (A)--(C) 是本文的具体语义结构论证；(D) 将其接入第~\ref{sec:compile}节的联合端口构造和标准知识编译接口。大小按表示项数计；序列化标识符时另计 $O(\log I)$ 位。$|X|$ 可以是 $\Theta(N^2)$，所以定理没有声称输出线性、求解器总内存恒定，或常数 $33k^2+k$ 足够小。

\subsection{证明：锁定块号而保留身份}
固定网络边 $u\to t$。标志 $f_u$ 只有一个写线程，因此其 CO 必须等于该线程的发布顺序，初始化在最前。将初始化记为块 0，阶段发布依次记块 $1,\ldots,n$。设 $t$ 的第 $i$ 次等待实际读取块 $q_i$。

RR coherence 使同线程有序读的来源写位置非降。相邻等待的观测位不同，而一个发布块内的值相同，因此 $q_1<\cdots<q_n$。最后一个等待读 $(n-1)\bmod2$，不可能来自值为 $n\bmod2$ 的块 $n$；于是 $q_n\le n-1$。在 $0,\ldots,n-1$ 中放置 $n$ 个严格递增整数，唯一可能是
\begin{equation}
 q_i=i-1\quad(1\le i\le n).\label{eq:block}
\end{equation}

旧核的非降来源链过滤可以得到这一结论：前向贪心计算每个读可接上合法前缀的最小写位置，后向贪心计算可接上后缀的最大位置；域内位于两界之间的候选能分别接上前、后缀。这是单链约束的精确过滤，不意味着该候选已经在全程序中可行。其他可靠规则可继续删域，却不能引入块外来源。剩余域若为空即可拒绝；否则块首发布是所有来源的反身 $H$ 前序，应用引理~\ref{lem:cover} 得到 (A)。来源块内的多个写从未合并。

\subsection{证明：沿网络传播共同因果}
取 $e$ 所在线程到 $f$ 所在线程的一条长为 $\ell\le d$ 的路径。设 $\sigma(e)=i$。阶段 $i$ 的任意事件，包括等待之后的数据操作，都 $\SB$-先于本线程阶段 $i+1$ 的首次发布。由 (A)，该发布先于第一个邻居阶段 $i+2$ 的等待。再沿该邻居程序序到阶段 $i+3$ 的首次发布，即可到第二个邻居阶段 $i+4$ 的等待。

每经过一条网络边推进两阶段，最终到达目标线程阶段 $i+2\ell$ 的等待，再沿程序序到其阶段 $i+2\ell+1$ 的全部事件。因此 $\sigma(f)\ge i+2d+1$ 时必有 $H(e,f)$。所需中间阶段都位于已知输入内；同线程情形直接由程序序得到。这证明 (B)。每跳使用两阶段是由等待在当轮最后发布之后决定的，不能倒走到已经发生的发布；本文不主张 $h$ 是最优界。

\subsection{证明：覆盖所有边种类的跨度}
令来源写 $w$ 位于阶段 $i$，读 $r$ 位于阶段 $j$。
若 $i-j>\delta$，(B) 给出 $H(r,w)$，来源应已删除。若 $j-i>\delta$，考虑长度恰为 $G$ 的整数窗口
\begin{equation}
 [j-h-G+1,\ j-h].\label{eq:pastwindow}
\end{equation}
它位于该地址的活动区间内。P3 给出其中一个同址写 $v$。两侧阶段距离均至少 $h$，故 $H(w,v)$ 和 $H(v,r)$，继而 $L(w,v)$；选择 $w$ 将形成 FR--HB 环。旧核删域规则因此排除这个过老来源。

对未决 CO，只要两个同址写的阶段差至少 $h$，(B) 就已将其方向加入 $L$，所以未决对不可能具有这种跨度。还须单独处理固定 CO 的生成边：若 $u,v$ 位于 $i<j$ 且 $j-i>\delta$，在
\begin{equation}
 [i+h,\ i+h+G-1]\label{eq:coverwindow}
\end{equation}
中由 P3 选同址写 $z$。两侧距离均至少 $h$，因此 $L(u,z)$ 和 $L(z,v)$；$(u,v)$ 不是 $L$ 的覆盖边。此步骤解释了为何写覆盖条件不能只由阶段密度取代。

相邻 $\SB$ 最多跨一阶段，因为每线程每阶段都有发布。初始化事件单独保留至多 $a$ 个；从初始化到各线程首事件的稀疏边足以保留其 HB 闭包。于是所有非初始化候选边跨度至多 $\delta$。

现在按组最大端点排序处理。设本组最大端点位于阶段 $j$。一个较早事件若仍活跃，就还关联某个未处理组。该组的最大端点不早于当前位置，而其跨度不超过 $\delta$，故这个活跃事件不早于阶段 $j-\delta$。尚未引入的未来事件不占当前袋。因此整个袋位于连续 $\delta+1$ 个阶段与初始化集合内，得到式~\eqref{eq:k}。这是实际处理顺序上的存活区间论证，而非只数某一条边的端点。至此证明 (C)；(D) 的构造和计费见下节。

\subsection{条件的边界与同一族上的结构差异}\label{sec:separation}
完整观测不能省略。若有三个发布块 $1,0,1$，却只有两个读观测 $0,1$，它们可以分别来自块 2、3；第一读不必来自初始化。该反例有经独立 RA 定义验证的完整赋值。写覆盖同样有实际作用：保留三线程 16 阶段的标志操作，但只在阶段 1、16 访问数据，则某些数据 $\Cov(L)$ 边跨度为 15。错误代入 $G=1$ 会得到 $\delta=11$；前提检查器拒绝此代入，使用真实覆盖间隔 $G=15$ 后结论成立。

为展示局部化确实改变候选结构，定义双线程族 $P(n,m)$，四个地址 $x,y,f_a,f_b$ 初值为 0，$m\ge2$。阶段 $i$ 令 $z_i=1$、$b_i=i\bmod2$，操作序列为
\begin{align*}
A: &W(x,z_i);\ W(f_a,b_i)_{0};\ W(y,z_i);\\
 &W(f_a,b_i)_{1};\ \cdots;\ W(f_a,b_i)_{m-1};\\
 &R(f_b,b_{i-1});\ R(x,z_i),\\
B: &W(y,z_i);\ W(f_b,b_i)_{0};\ W(x,z_i);\\
 &W(f_b,b_i)_{1};\ \cdots;\ W(f_b,b_i)_{m-1};\\
 &R(f_a,b_{i-1});\ R(y,z_i).
\end{align*}
所有 $W$ 为 release，所有 $R$ 为 acquire。将 $z_i$ 改为 $i$ 得版本载荷族 $\pv(n,m)$。它使数据读仅匹配同轮两位写者，但不区分这两位写者；标志仍只有交替位。版本值需要 $O(\log n)$ 位，所以该变体不保持固定有限的数据值域。

\begin{cprop}[宽度差异与多种支持]\label{prop:separation}
对 $P(n,m)$ 和 $\pv(n,m)$，$K_0$ 的候选图包含 $K_{n,n}$，其最大事件接口随 $n$ 无界。$\ka$ 则满足定理~\ref{thm:main}，其中 $d=1,G=1,b=2(m+4),a=4$，故 $k\le16m+68$。两族至少具有 $4^n m^{2(n-1)}$ 个合法原始赋值；每个非首轮等待的 $m$ 个上轮来源均有合法支持。
\end{cprop}
\begin{cproof}
$K_0$ 锁定标志来源块后，非首轮等待仍有 $m$ 个来源；数据读至少保留本线程本轮与另一线程本轮两个来源。因此除初始化外，不产生跨线程必选 RF。$H$ 没有非初始化跨线程边，$L$ 只固定线程内写序。在地址 $x$ 上，A 的 $n$ 个写与 B 的 $n$ 个写之间 $n^2$ 个 CO 比较全部未决，形成 $K_{n,n}$。任意事件排列中，某一侧第一次全部引入、另一侧尚未全部引入的切口上，已完成侧的 $n$ 个点均有未来邻居，不能退休。局部化界直接由 P1--P3 得到。

对模型数下界，令每个数据读选择本线程本轮的第一项写，等待独立选择任一上轮发布，首轮读初始化。两个数据地址的 CO 均按轮拼接，同轮两位写者的先后各自任选；标志 CO 按程序序。HB 不倒退轮号，跨线程非初始化 HB 至少推进一轮，同轮数据写之间没有跨线程 HB。由此 WW、WR、RW、RR coherence 均成立，详细分类见附录~\ref{app:models}。每轮两个数据地址给出四种独立 CO 选择，非首轮两次等待各有 $m$ 种来源，得到所列下界。
\end{cproof}

该差异仅针对明确定义的 $K_0$ 候选表示；它不是 SAT、BDD、一般 may/must 分析或所有布局编译器的复杂度下界。两族也确实体现弱内存：由 (A)，首轮等待全部读取初始化。对同步边 $u\to t$，任何 SC 全序必须有“$t$ 首次发布 $<$ $t$ 读取 $u$ 初值 $<$ $u$ 首次发布”；沿有向环组合得到矛盾。因此满足 P1、P2 的这些完整观测在 SC 下不可达。上述合法赋值则证明所选数据观测在 RA 下可达。
